%Locate this file one level lower than the folder with references:
%		-https://svn.oss.deltares.nl/repos/openearthtools/trunk/references
%which must be checkout with in a folder named <00_references>.
%E.g., locate this file renamed as <memo.tex> in a structure:
%		-<#/00_references/>
%		-<#/01_memo/memo.tex>
%
%Folder <#00_reference/02_textree> must be in the MikTeX path. See <#00_reference/02_textree/readme>
%
%
%
%
%
%
%----------------------------
%---------DOCUMENT TYPE
%----------------------------
%

\newcommand{\isreport}{1} %0=memo; 1=report
\newcommand{\addAppendix}{1} %0=no appendix; 1=add appendix
\newcommand{\addCitation}{0} %0=no; 1=yes
\newcommand{\addmemosignatures}{0} %0=no; 1=yes

%Add colon in \mySubtitle if desired or leave empty if no substitle is desired. 
\def\myTitle{Verkenning tool voor effecten ijsvorming en dwarsstroming}
\def\mySubtitle{}

\def\myOrganisationi{Deltares}
\def\myAuthori{Robert Groenewege}
\def\myAuthorii{}
\def\myPhonei{+31\,(0)64\,691\,1978}
\def\myEmaili{robert.groenewege@deltares.nl}

\def\myVersion{0.1}
\def\myDate{\today}

%
%
%
%----------------------------
%---------DOCUMENT CLASS
%----------------------------
%

\ifnum \isreport=1
	\documentclass[dutch,nosummary,signature]{deltares_report}
	%language: english, ducth, spanish
	%nosummary: writes no executive summary 
	%signature: adds box for signatures
\else
	\documentclass{deltares_memo}
\fi

%
%
%
%----------------------------
%---------PACKAGES AND COMMANDS
%----------------------------
%

\input{00_references/abbreviations}
\input{00_references/packages}
\usepackage[colorinlistoftodos]{todonotes}
\usepackage{pdfpages}
\usepackage{tablefootnote}

%
%
%
%----------------------------
%---------DOCUEMENT INFORMATION
%----------------------------
%

%%%%%%%%
%---------REPORT INPUT
%%%%%%%%
%

\ifnum \isreport=1

\def\authorList{\myAuthori{}, \myAuthorii{}} %Where does this go in the layout?

%adding a different figure on the cover
%
%\renewcommand{\FrontCover}{\includegraphics[height=182mm,width=182mm]{../00_figures/12_other/DJI_0017.JPG}}
%\renewcommand{\FrontCircle}{\includegraphics[height=182mm,width=182mm]{../00_references/01_layouts/cover/cover_transparant_circle.pdf}}		
%\renewcommand{\FrontCover}{\includegraphics[height=182mm,width=182mm]{pictures/Deltares_rapport_omslag_gennepperhuis.jpg}}															   
%prevent sentence break between pages		
\widowpenalties 1 10000
\raggedbottom
	
\begin{document}

\title{\myTitle}
\subtitle{\mySubtitle}
\author{\myAuthori{} \\
			\myAuthorii{}}
\partner{}
\coverPhoto{}
\date{\today}
\version{0.1}

%\foreach \var [evaluate=\var as \myindex using {int(\var-1)}] in {1} {
	%\pgfmathparse{\authorList[\myindex]}\pgfmathresult
%}

\authori{\myAuthori{}}
\organisationi{\myOrganisationi}
%\authorii{\myAuthorii{}}
%\organisationii{\myOrganisationi}

%revision table
\revieweri{}
\datei{\myDate}
\versioni{0.1} 
\approvali{}  
\publisheri{}
\authorii{\myAuthorii{}}

\client{Rijkswaterstaat}
\contact{Arjan Sieben}
\keywords{}
\reference{}
\classification{}
\status{concept}
\disclaimer{}
%\summary{} 
 
\documentid{}
\projectnumber{11211565-010}


%%%%%%%%
%---------MEMO INPUT
%%%%%%%%
%


\else

\begin{document}
%\memoTo{Aukje Spruyt, Johan Boon}
%\memoConfidentialUntil{}
%\memoName{Memo} %Name in header. If commented out, default "Memo" is used. 
%\memoDate{\today\currenttime}
%\memoVersion{\myVersion{}} %11206793-013-ZWS-0001_v0.1-groynes.docx
%\memoFrom{\parbox[t]{3cm}{
\myAuthori{} 
\myAuthorii{} \\ 
}}
\memoTelephone{\parbox[t]{3cm}{
\myPhonei{} 
%\myPhoneii{} \\
}}
\memoEmail{\parbox[t]{3cm}{
\myEmaili{} 
%\myEmailii{}
}}
\memoSubject{\myTitle{}\mySubtitle{}}
\memoCopy{}

\fi

%
%
%
%----------------------------
%---------COVER
%----------------------------
%
%\svnInfo $Id: report.tex 83 2022-12-05 10:24:56Z ottevan $
\deltarestitle
\ifnum \isreport=1
	\newpage
\fi
\ifnum \isreport=0
\ifnum \addmemosignatures=1
\begin{tabular}{p{\textwidth/8}|p{0.175\textwidth}|p{0.2\textwidth}|p{0.2\textwidth}|p{0.2\textwidth}}
    \rowcolor{dblue1}  \textbf{Document version} & \textbf{Date} & \textbf{Author} & \textbf{Reviewer} & \textbf{Approval} \\
    \topline
    0.1   & & Robert Groenewege   &   &  \\
		\midline
		& & & \\
    \midline                                                 
\end{tabular}
\fi
\fi
%
%
%
%----------------------------
%---------START DOCUMENT
%----------------------------
%

\def\RijnFigDir{../examples/c01 - Waal/figures/}
\def\MaasFigDir{../examples/c02 - Maas/figures/}
\def\RMMFigDir{../examples/c03 - RMM/figures/}
\def\NVOMaasFigDir{../examples/c04 - NVO Maas/figures/}

%Second input is the level of the section.
\gensection{\isreport}{1}{Projectomschrijving}
\label{sec:projectplan}

\gensection{\isreport}{2}{Aanleiding}
Dit project kent twee aanleidingen:
\begin{itemize}
	\item de beoordeling van de invloed van ingrepen op ijsafvoer blijkt in de praktijk lastig uit te voeren, een eenduidige, efficiënte methodiek ontbreekt. 
	\item voor beoordeling van dwarsstroomeffecten zijn (en worden) door initiatiefnemers verschillende aanpakken ontwikkeld. Met een uniforme aanpak kan de kwaliteit beter worden geborgd. 
\end{itemize}

In 2024 is door KRW projecten voor de Rijntakken een aanpak gevolgd voor het kwantitatieve onderdeel van de beoordeling van effecten op ijsafvoer. Voor SITO Rivierkunde is de invloed van onregelmatige oevers op rivierfuncties (veilige afvoer van water, sediment en ijs, dwarsstroming, vlot en veilig varen en ondersteuning laagwaterstanden) onderzocht \citep{Groenewege25}. De beide ontwikkelingen geven perspectief op generiek gebruik dat in een rekentool geformaliseerd kan worden.

\gensection{\isreport}{2}{Doel}
Doel van dit deelproject is het ontwikkelen van een eenduidige aanpak met rekentool, voor de RBK-bepaling van dwarsstroom- en ijsafvoereffecten, waarmee naar gelang de behoefte, eenduidig 
\begin{itemize}
	\item op uniforme wijze de grootte van dwarsstroming en dwarsstromingseffecten kan worden bepaald conform de specificaties in het RBK; 
	\item de relevante variabelen voor het schatten van de invloed van maatregelen op de doorvoer van ijs kunnen worden bepaald, conform de specificaties in het RBK, inclusief correcties voor invloeden van benedenstrooms ijsdek en lokale bodemveranderingen.
\end{itemize}

Dit is een gecombineerde aanpak van twee verschillende aspecten (ijsafvoer en dwarsstroming op de vaarweg) vanwege een redelijke overlap in de voor evaluatie relevante variabelen.

Bij de rekentool gaat het in 2025 om een serie van eenvoudige scripts, en nadrukkelijk nog niet om een tool dat uitgeleverd kan worden voor gebruik door ingenieursbureaus. Mits de resultaten veelbelovend zijn kan in een volgende stap erover worden nagedacht hoe dit omgezet kan worden in een tool dat door de markt toegepast kan worden, bijvoorbeeld een aanvulling van het set tools onder D-FAST. Daarbij moet dan ook worden bekeken hoe het beheer en onderhoud van de tool kan worden gewaarborgd. Wel is in 2025 geprobeerd om al zo veel mogelijk aan te sluiten bij een implementatie in de D-FAST productlijn.

\gensection{\isreport}{2}{Werkzaamheden}
\label{}

De volgende werkzaamheden zijn uitgevoerd:
\begin{itemize}
	\item Definitie van de tool ten aanzien van de inhoud (toepassingsgebied, beperkingen, enz.) en het gebruik (processing D-HYDRO-rekenresultaten, format rapportage)
	\item Scripting en test
	\item Toepassing voor de Rijn, Maas en Rijn-Maasmonding: rapportage en interpretatie van rivierstukken middels eerste kwalificatie
	\item Rapportage met een beschrijving van de tool en de resultaten van de toepassing
\end{itemize}

Dit leidt tot de volgende opgeleverde producten:
\begin{itemize}
	\item prototype van de tool
	\item rapport (voorliggend document)
	\item instructies voor installeren en gebruik (zie User Manual)
\end{itemize}
\textit{Procesafspraak}: De RWS klankbordgroep wordt in alle fases van het deelproject betrokken.

\gensection{\isreport}{1}{Achtergrond}
\label{}

\gensection{\isreport}{2}{Afvoer van ijs}
\label{achtergrond_ijs}
Het RBK stelt dat een goede geleiding van ijs en water gewaarborgd moet blijven (RBK, sectie 1.5):

\begin{quote}
"De volgende ontwerpprincipes zijn relevant voor een goede afvoer van ijs:
	\begin{itemize}
		\item In het stroomvoerend profiel mag de ingreep voor afvoeren vanaf bankfull tot grofweg 75 jaar herhalingstijd (Lobith van 4000 tot 8000 m3/s, Borgharen van 1500 tot 2800 m3/s), ook in scenario’s met benedenstrooms ijsdek, de Froude getallen niet verlagen tot onder 0.08, om de kans op ontwikkeling van ijsdammen niet te verhogen); 
		\item Verander lokaal de normaalbreedte van de rivier niet. Lokale versmallingen kunnen de ijsafvoer blokkeren;
		\item Laat de gestrekte oevers in stand, met name in het splitsingspuntengebied;
		\item Voorkom dat grote stukken ijs massaal vanuit de nevengeul in de hoofdgeul kunnen stromen en op die manier blokkades gaan vormen;
		\item Voorkom de vorming van ondieptes in het zomerbed."
	\end{itemize}
\end{quote}

De beoordeling van de doorvoer van ijs is in de meeste gevallen grotendeels kwalitatief. Echter, twee ontwerpprincipes uit het Rivierkundig Beoordelingskader (versie 6.0) kunnen voor zover weergegeven in D-HYDRO, eenduidig kwantitatief worden toegepast:
\begin{enumerate}
	\item In het stroomvoerend profiel mag de ingreep voor afvoeren vanaf bankfull tot grofweg 75 jaar herhalingstijd (Lobith van 4000 tot 8000 m3/s), ook in scenario’s met benedenstrooms ijsdek, de Froude getallen niet verlagen tot onder 0.08, om de kans op ontwikkeling van ijsdammen niet te verhogen.
	\item Verander lokaal de normaalbreedte van de rivier niet. Lokale versmallingen kunnen de ijsafvoer blokkeren. Laat de gestrekte oevers in stand, met name in het splitsingspuntengebied.
\end{enumerate}

Ad 1) Het Froude-getal karakteriseert de relatieve convectie van stroming met ijs. Bij stroming met ijs tegen ijsdekken kunnen volgens het criterium van Kivisild ijsdammen verwacht worden voor Froude-getallen onder 0.015 á 0.150 \citep{Termes91}. Het gemiddelde van deze range, 0.08, lijkt een goede grenswaarde voor ijsdamvorming \citep{Zagonjolli19}.

Dit betreft de overgang van stroming met los ijs naar stroming onder een vast ijsdek. Een haperende doorvoer van ijs leidt tot vorming van ijsdammen. De vrije stroming ondergaat opstuwing vanuit het benedenstrooms ijsdek en de invloed van de bodemveranderingen als gevolg van de ingreep. Het Froude-getal dat stroming karakteriseert met D-HYDRO resultaten zonder ijsdek en bodemverandering heeft dus twee correcties nodig. Deze zijn beschreven in bijlage \ref{app:ijscorrecties}. 

Ad 2) De ijsafvoer stagneert bij een afnemend ijs-transporterend vermogen. In trajecten met getij is dit het geval gedurende het wantij. In stationaire stroming is dat ter plekke van verbredingen en splitsingen, bij obstakels, op ondiepten en in sterk gekromde stroming. Een constante normaalbreedte met vloeiend verlopende normaallijnen borgt een goede ijsafvoer, omdat dan gradiënten in stroming en daarmee de gradiënten in ijs-transporterend vermogen beperkt zijn.

Vanwege de invloed op meerdere takken weegt handhaving van deze normaalbreedten (en de normaallijnen die deze lijnen definiëren) in en rondom de splitsingspuntgebieden zeer zwaar. Daarbuiten kan, om de continuïteit in ijs-transporterend vermogen te borgen, de invloed van maatregelen op de ijsafvoer van stationaire stroming (dus zonder getij) voldoende gekarakteriseerd worden met de gradiënt in grootte en richting van de stroomsnelheid. Een ontwerp mag in het rivierstuk van de ingreep, bij de genoemde rivierafvoeren niet leiden tot meer of grotere van deze gradiënten.

\gensection{\isreport}{2}{Dwarsstroming}
Het kwantitatieve onderdeel van de beoordeling van invloeden op de vaarweg is veelal gericht op bepaling van dwarsstroomsnelheden. Immers, van ingrepen wordt verwacht dat deze bij dwarsstroomdebieten groter dan 50 m³/s, de dwarsstroming in de vaarweg niet boven de richtlijn van 0.15 m/s verhogen, tenzij hierdoor de padbreedte van passerende schepen niet meer dan een halve scheepsbreedte toeneemt. Bij debieten kleiner dan 50 m³/s is dit een richtlijn van 0.30 m/s. Dit zijn de criteria in RBK6, toegespitst op nevengeulen e.d. In RVW zijn criteria opgenomen voor geconcentreerde dwarsstroming: kleinere debieten met grotere uitstroomsnelheid (bv. in/uitlaat koelwater bij een centrale gemaal/ riooloverstort/ inlaatpunt, e.d.). Binnen WVL loopt er momenteel een actie om deze de criteria vanuit deze twee invalshoeken beter op elkaar aan te laten sluiten (pers. comm. klankbordgroep). De tool moet dan ook anticiperen op aanpassing van criterium.

De toepassing voor RBK beoordeling begint veelal met het in kaart brengen van stroombeelden en stroomsnelheden voor de situatie zonder ingrepen en voor de situatie met ingrepen. De dwarsstroomsnelheid op de rand van de vaarweg (zoals gedfinieerd in het RBK) moet voor een aantal kenmerkende rivierafvoeren en getijverlopen (Rijn-Maasmonding) worden gepresenteerd in grafieken. Deze grafieken geven de rivierbeheerder inzicht in de effecten van de ingreep op dwarsstromingen. Ook de representatieve dwarsstroming wordt door de tool berekend en gepresenteerd.

\gensection{\isreport}{1}{Beschrijving van de rekentool}
In dit hoofdstuk wordt beschreven wat de beoogde werking en functionaliteit van de rekentool is.

\gensection{\isreport}{2}{Inleiding}
De broncode van de tool is geschreven in de programmeertaal Python. Dit maakt het makkelijker om in een later stadium de tool binnen de D-FAST familie, ook geschreven in Python, te implementeren. Om deze reden wordt er ook zoveel mogelijk gebruik gemaakt van al bestaande functies van D-FAST. 

De tool is te gebruiken als "command line interface" (CLI) zonder een grafische gebruikersomgeving (GUI). De tool wordt meegeleverd met instructies voor installeren en gebruik. De benodigde invoer van de tool betreft het volgende, vergelijkbaar met D-FAST MI:
\begin{itemize}
	\item map- of fourier uitvoerbestanden van D-HYDRO simulaties;
	\item configuratiebestand van de riviertakken (naam, trajecten, afmetingen van schepen, enz.);
	\item configuratiebestand van de analyse (hiermee kunnen gebruikers verwijzen naar de relevante D-HYDRO simulaties).
\end{itemize}

De uitvoer van de tool bestaat uit figuren en netCDF bestanden. In de volgende secties is dit verder uitgewerkt.

\gensection{\isreport}{3}{Software vereisten}

Uit de feedback vanuit de RWS klankbordgroep zijn enkele algemene vereisten aan de tool gedefinieerd:
\begin{enumerate}
	\item Geef gebruikers de mogelijkheid om de code aan te passen/suggesties te doen voor verbetering. Beheerder kan ze dan al/of niet opnemen in de officiële versie.
	\item Formuleringen en criteria moeten 1 op 1 overeenkomen met RBK.
	\item De tool moet anticiperen op aanpassing van criteria voor de dwarsstroming.
	\item De tool moet geschikt zijn voor het berekenen van de representatieve dwarsstroming.
	\item Voor dwarsstroming moet de tool ook een verschilplot (bijvoorbeeld op de secundaire as) genereren.
	\item De tool moet geen afstand over de lijn visualiseren maar rivierkilometer (rkm) op de x-as.
	\item De tool moet de mogelijkheid bieden om modelresultaten langs de raai-km van de rivier te plotten.
	\item De tool moet de mogelijkheid bieden voor een inverse rkm-as.
	\item De gebruiker moet zelf lijnen kunnen definiëren, waarlangs de modelresultaten worden geplot.
	\item Voor de Rijn-Maasmonding moet de tool de maxima kunnen berekenen voor eb en vloed.
\end{enumerate}


\gensection{\isreport}{4}{Afvoer van ijs}
Om de afvoer van ijs te beoordelen, kunnen de amplitude en richting van stroomsnelheden ten eerste gevisualiseerd worden langs drie lijnen (ééndimensionaal):
\begin{itemize}
	\item de twee lijnen die al gebruikt worden voor het in beeld brengen van dwarsstroming (of, als deze ontbreken, de waterlijn bij bankvullende afvoer). 
	\item de lijn over de rivieras die wordt gebruikt voor de waterstandseffecten op de as. 
\end{itemize}
Stroomsnelheden langs deze lijnen worden door de tool over voldoende lengte gevisualiseerd om een goede vergelijking van de waarden ter plekke van de ingreep met waarden op ongestoorde oeverdelen mogelijk te maken. 

Ten tweede produceert de tool figuren van Froude getallen voor de situatie zonder ingreep en voor de situatie met ingreep, met een classificatie van de toe- en afname in verschillende klassen (Figuur \ref{fig:Froude}). Dit is inclusief de twee benodigde correcties (zie \ref{app:ijscorrecties}).

\begin{figure*}[hbt!]
\centering
\includegraphics[width=\textwidth]{01_figures/HKV_Froude.png}
\captionsetup{justification=centering}
\caption{Beoogde visualisatie van Froude getallen. Bron: HKV}
\label{fig:Froude}
\end{figure*}

\gensection{\isreport}{4}{Dwarsstroming}
De beoordeling van dwarsstroming is uitvoerig beschreven in bijlage 7 van het RBK. De tool automatiseert de stappen voor de bepaling en presentatie van dwarsstroming. De geproduceerde figuren volgen de opzet van Figuur \ref{fig:dwarsstroming}, maar visualiseren óók het variërende criterium (0.15 of 0.3 $m/s$) dat hoort bij het lokale dwarsstroomdebiet. De dwarsstroming wordt gepresenteerd langs de lijn(en) die de gebruiker zelf opgeeft (zie volgende alinea). Overschrijdingen van het (lokale) criterium worden ook apart weggeschreven in een Excel bestand. Het is mogelijk om in één figuur zowel de lijn voor een situatie zonder ingreep als de lijn voor een situatie met ingreep te visualiseren, zodat het effect van de ingreep inzichtelijk is. Op de secundaire as wordt het verschil tussen de situatie zonder en met ingreep toegevoegd. Dit leidt tot drie lijnen:
\begin{enumerate}
	\item Referentie (zwarte lijn)
	\item Met ingreep (blauwe lijn)
	\item Verschil = ingreep - referentie (rode lijn)
\end{enumerate}

Bij RWS-ZN/Maas is het normprofiel bepalend voor de RBK-toets op dwarsstroming. Bij RWS-ON/Rijntakken is de rand van de vaarweg de norm. Deze wordt gemarkeerd door de bakenlijn langs de kribbakens. Bij de Waal en Nederrijn/Lek komt de bakenlijn vrijwel overeen met de normaallijn. Bij de IJssel ligt de normaallijn enkele meters uit de bakenlijn en zijn er ook trajecten waar de bakenlijn en/of de normaallijn is gewijzigd. Vanwege de diversiteit aan mogelijke profiellijnen, is ervoor gekozen om in de tool de gebruiker de mogelijkheid te bieden zelf één of meerdere lijnen aan te leveren.

\begin{figure*}[hbt!]
\centering
\includegraphics[width=\textwidth]{01_figures/RBK_dwarsstroming.png}
\captionsetup{justification=centering}
\caption{Voorbeelden van de visualisatie van dwarsstroming. Boven: bron; RBK 6.0, Bijlage 7, Figuur 1}
\label{fig:dwarsstroming}
\end{figure*}

\gensection{\isreport}{4}{Toepassingsgebied}
In principe kan de tool uitgevoerd worden voor alle riviertakken waarvoor maatgevende schipafmetingen zijn gedefinieerd (in het RBK). Dit betreft:

\begin{itemize}
	\item Bovenrijn (Rkm 859-867)
	\item Waal (Rkm 868-951)
	\item Pannerdensch-Kanaal (Rkm 868-879)
	\item Nederrijn-Lek (Rkm 880-989)
	\item IJssel (Rkm 880-1000)
	\item Zwarte Water
	\item Merwedes (Rkm 951-980)
	\item Maas (Rkm 16-227)
	\item Amer 
	\item Haringvliet 
	\item Hollands Diep
\end{itemize}

Echter, in de praktijk ondersteunt de tool op dit moment alleen riviertakken waar de invloed van getij verwaarloosbaar klein is ten opzichte van rivierafvoer. Anticiperend op implementatie van de tool in D-FAST, is er in deze studie nog weinig aandacht besteed aan het uitbreiden van dit toepassingsgebied. Voor D-FAST-MI wordt hier namelijk al aan gewerkt, maar ten tijde van schrijven was deze functionaliteit nog niet officieel gereed. Op termijn is het de bedoeling om hierop aan te sluiten.

\gensection{\isreport}{4}{Beperkingen}
De volgende beperkingen zijn bekend:
\begin{itemize}
	\item Als de analyse wordt uitgevoerd op een gebied met meerdere rivierassen, zoals in de RMM, mogen er geen dubbelingen in de rivierkilometers voorkomen. Er moet dus gekozen worden voor één as(lijn).
	\item Voorlopig is het nog niet mogelijk om met de tool percentielwaarden uit te rekenen. Als alternatief kan de gebruiker bijvoorbeeld zelf de laatste tijdstap van het map.nc bestand overschrijven met de gewenste waarden; hiervoor wordt de QGIS Crayfish plugin aangeraden.
\end{itemize}
	
\gensection{\isreport}{1}{Validatie}
In dit hoofdstuk worden voorbeelden gegeven waarmee de visualisatie van resultaten wordt toegelicht en de werking van de tool is gevalideerd. Ten eerste wordt de hypothetische natuurvriendelijke oever van \citep{Groenewege25} in de Maas bij rivierkilometer 188 gepresenteerd. Ten tweede is er een eerste karakterisering gemaakt van dwarsstroming en Froude-getallen van de gehele Rijn, Maas en Rijn-Maasmonding. De exacte visualisatie van resultaten door de tool kan mogelijk nog wijzigen ten opzichte van het eindproduct.

\gensection{\isreport}{2}{Natuurvriendelijke oever Maas}
De natuurvriendelijke oever (NVO) betreft een symmetrische erosiekom van 50 m breed en 100 m lang. De resultaten worden ter illustratie voor één afvoersom gepresenteerd (S2100); overige afvoersommen kunnen in \citep{Groenewege25} worden gevonden. 

\gensection{\isreport}{3}{Afvoer van ijs}
In deze sectie worden de effecten van de NVO op de afvoer van ijs getoond.

\gensection{\isreport}{4}{1D profielen}

In Figuur \ref{fig:NVO_profiel0} en Figuur \ref{fig:NVO_profiel1} worden de randen van de vaarweg getoond, waarlangs stroomsnelheid en -richting worden gepresenteerd. De bodemligging in de plansituatie fungeert als achtergrond.

\begin{figure*}[hbt!]
\centering
\includegraphics[width=0.5\textwidth]{\NVOMaasFigDir/profile0_location.png}
\captionsetup{justification=centering}
\caption{Locatie van profiel 0}
\label{fig:NVO_profiel0}
\end{figure*}

\begin{figure*}[hbt!]
\centering
\includegraphics[width=0.5\textwidth]{\NVOMaasFigDir/profile1_location.png}
\captionsetup{justification=centering}
\caption{Locatie van profiel 1}
\label{fig:NVO_profiel1}
\end{figure*}

Figuur \ref{fig:NVO_stroming_p0} en Figuur \ref{fig:NVO_stroming_p1} tonen stroomsnelheid en -richting langs de profielen (zwarte en blauwe lijnen). De x-as is omgekeerd (stroming van rechts naar links). Stroomsnelheid begint bij 0, en stroomrichting varieert van -90 tot 90 graden (loodrecht) ten opzichte van de profiellijn\footnote{Er is voor dit bereik gekozen vanwege eb en vloed; in de broncode zit geen afhankelijkheid van de oriëntatie van de profiellijn (stroomopwaarts of -afwaarts). Om te beoordelen of de stroming omkeert door de ingreep (richtingsverandering van meer dan 90 graden), moeten aanvullend de 2D modelresultaten worden bekeken.}. De situatie zonder ingreep wordt aangeduid met 'Referentie' en de situatie met ingreep met 'Plansituatie'. Het verschil hiertussen wordt in rood rechts op een secundaire y-as weergegeven, die gecentreerd is rond 0. Met deze figuren wordt duidelijk dat snelheidsgradiënten worden versterkt ter plekke van de NVO op rkm 188. Aan de kant van de NVO zien we tevens verandering van stromingsrichting van meer dan 10 graden.

\begin{figure*}[hbt!]
\centering
\includegraphics[width=0.5\textwidth]{\NVOMaasFigDir/C2_profile0_velocity_angle.png}
\captionsetup{justification=centering}
\caption{Stroomsnelheid en -richting langs profiel 0 (S2100)}
\label{fig:NVO_stroming_p0}
\end{figure*}

\begin{figure*}[hbt!]
\centering
\includegraphics[width=0.5\textwidth]{\NVOMaasFigDir/C2_profile1_velocity_angle.png}
\captionsetup{justification=centering}
\caption{Stroomsnelheid en -richting langs profiel 1 (S2100)}
\label{fig:NVO_stroming_p1}
\end{figure*}

\gensection{\isreport}{4}{Froude getallen}
Figuur \ref{fig:NVO_Froude} laat de Froude getallen zien rondom de NVO, zonder correcties voor waterstandsopstuwing door een ijsdek of voor lokale bodemveranderingen. Het verschil tussen de plansituatie en de referentie wordt getoond in Figuur \ref{fig:NVO_Froude_diff1}. Binnen de erosiekom zien we een toename van Froude getallen van < 0.08 naar >= 0.08. Hier direct benedenstrooms van worden Froude getallen verlaagd van > 0.08 naar <= 0.08. Het toevoegen van de correctie voor waterstandsopstuwing (\ref{fig:NVO_Froude_diff2}) verkleint de verschillen met 30\% (zie Appendix \ref{app:ijscorrecties}) en vervolgens levert de correctie voor lokale bodemverandering (\ref{fig:NVO_Froude_diff3}) in dit geval relatief weinig op.

\begin{figure*}[!h]
\centering
\includegraphics[width=0.5\textwidth]{\NVOMaasFigDir/C2_intervention_Froude.png}
\captionsetup{justification=centering}
\caption{Froude getallen in plansituatie, zonder correcties (S2100)}
\label{fig:NVO_Froude}
\end{figure*}

\begin{figure*}[hbt!]
\centering
\includegraphics[width=0.5\textwidth]{\NVOMaasFigDir/C2_difference_Froude.png}
\captionsetup{justification=centering}
\caption{Verschil in Froude getallen tussen plansituatie en referentie, zonder correcties (S2100)}
\label{fig:NVO_Froude_diff1}
\end{figure*}

\begin{figure*}[hbt!]
\centering
\includegraphics[width=0.5\textwidth]{\NVOMaasFigDir/C2_difference_Froude_wateruplift.png}
\captionsetup{justification=centering}
\caption{Verschil in Froude getallen tussen plansituatie en referentie, met correctie voor wateropstuwing door benedenstrooms ijsdek (S2100)}
\label{fig:NVO_Froude_diff2}
\end{figure*}

\begin{figure*}[hbt!]
\centering
\includegraphics[width=0.5\textwidth]{\NVOMaasFigDir/C2_difference_Froude_wateruplift_bedchange.png}
\captionsetup{justification=centering}
\caption{Verschil in Froude getallen tussen plansituatie en referentie, met alle correcties (S2100)}
\label{fig:NVO_Froude_diff3}
\end{figure*}

\gensection{\isreport}{3}{Dwarsstroming}
In deze sectie worden de effecten van de NVO op de dwarsstroming getoond.

\begin{figure*}[hbt!]
\centering
\includegraphics[width=\textwidth]{\NVOMaasFigDir/C2_profile0_transverse_flow.png}
\captionsetup{justification=centering}
\caption{Dwarsstroming langs profiel 0 (S2100)}
\label{fig:NVO_Froude_ref}
\end{figure*}

\begin{figure*}[hbt!]
\centering
\includegraphics[width=\textwidth]{\NVOMaasFigDir/C2_profile1_transverse_flow.png}
\captionsetup{justification=centering}
\caption{Dwarsstroming langs profiel 1 (S2100)}
\label{fig:NVO_Froude_diff}
\end{figure*}

\gensection{\isreport}{1}{Karakterisering van de Rijn}
In dit hoofdstuk wordt een eerste karakterisering gemaakt van stroming in de gehele Rijn, Maas en Rijn-Maasmonding met betrekking tot de afvoer van ijs en dwarsstroming. Voor alle 

\gensection{\isreport}{2}{Opzet}
Voor de berekeningen is gebruik gemaakt van het \texttt{dflowfm2d-rijn-j24\_6-v1a} model. De afvoersommen S4000, S6000 en S8000 zijn doorgerekend om het bereik "Lobith van 4000 tot 8000 m3/s" (zie sectie \ref{achtergrond_ijs}) te representeren. 

\gensection{\isreport}{2}{Resultaten}

\gensection{\isreport}{3}{Afvoer van ijs}

\gensection{\isreport}{4}{1D profielen}

\gensection{\isreport}{4}{Froude getallen}

\gensection{\isreport}{43}{Dwarsstroming}

\gensection{\isreport}{2}{Maas}

\gensection{\isreport}{3}{Opzet}
Voor de berekeningen is gebruik gemaakt van het \texttt{dflowfm2d-maas-beno2\2_6-v2a} model. De afvoersommen S1700, S2100 en S2500 zijn doorgerekend om het bereik "Borgharen van 1500 tot 2800 m3/s" (zie sectie \ref{achtergrond_ijs}) te representeren. 

De randen van de vaarwegen zijn 

\gensection{\isreport}{3}{Resultaten}

\gensection{\isreport}{4}{Afvoer van ijs}

\gensection{\isreport}{5}{1D profielen}

\gensection{\isreport}{5}{Froude getallen}

\gensection{\isreport}{4}{Dwarsstroming}

%\gensection{\isreport}{2}{Rijn-Maasmonding}
%
%\gensection{\isreport}{3}{Opzet}
%Voor de berekeningen is gebruik gemaakt van het \texttt{dflowfm2d-rmm\_vzm-j17\_6-v2a} model. Er is gerekend met gemiddeld getij, en Rijnafvoer bij Tiel van 4062.8, 5417.3 en 6616.9 m3/s en Maasafvoer bij Lith van 1235, 1742 en 2248 m3/s, om het bereik "Lobith van 4000 tot 8000 m3/s" en "Borgharen van 1500 tot 2800 m3/s" (zie sectie \ref{achtergrond_ijs}) te representeren. Het Fourier bestand is enigszins aangepast om de juiste uitvoer te krijgen, en 1 verouderd keyword ("transportmethod") is in de MDU aangepast. In het RBK staan alleen voor de takken Haringvliet, Hollands Diep, Amer en de Merwedes maatgevende schipafmetingen genoemd; zodoende is de analyse alleen voor deze takken uitgevoerd. De analyse beperkt zich hier nog tot Froude getallen in 2D.
%
%\gensection{\isreport}{3}{Resultaten}
%
%\gensection{\isreport}{4}{Afvoer van ijs}
%
%\gensection{\isreport}{5}{Froude getallen}


%
%
%
%----------------------------
%---------GLOSSARY
%----------------------------
%
%\section{List of Mathematical Symbols}
%\printglossary[title=]

%define
% \newglossaryentry{H}{
  % name={$H$},
  % description={weir height. Distance between the crest and the bed level on the downstream side.},
  % sort={H}
% }

%use
% \gls{H}

%
%
%
%----------------------------
%---------BIBLIOGRAPHY
%----------------------------
%
\FloatBarrier
\gensection{\isreport}{1}{Referenties}
\DeclareRobustCommand{\van}[3]{#3}
\bibliography{00_references/references}


\vfill
\ifnum \addCitation=1
Please cite this work using: \\
\\
\texttt{@TechReport\{[Add\_bibtex\_reference],\\
~author~~~= \{\myAuthori{} and \myAuthorii{}\},\\
~title~~~=  \{\myTitle{}: \mySubtitle{}\},\\
~number~~~=      \{\myVersion{}\},\\
~year~~~~=        \{\myDate{}\},\\
~institution= \{{Deltares, Delft, the Netherlands}\}\\
\ifnum \isreport=1
~type~~~~=  \{\{T\}ech. \{R\}ep.\},\\
\else
~type~~~~= \{\{M\}emo\}, \\
\fi
~\}
}
\fi

%
%
%
%----------------------------
%---------APPENDIX
%----------------------------
%
\ifnum \addAppendix=1

\newpage
\appendix
\renewcommand\thesection{\AlphAlph{\value{section}}}

\gensection{\isreport}{1}{Correcties van modelresultaten voor de afvoer van ijs}
\label{app:ijscorrecties}

\textbf{Correctie 1}\\
De verandering in waterdiepte $h$ [$m$] door opstuwing van een benedenstrooms ijsdek wordt naar \citet{Zagonjolli19} pragmatisch geschat met 
\begin{equation}
	{\frac{h_{\text{met ijsdek}}}{h_{\text{zonder ijsdek}}}} \approx \left(\frac{C_{\text{zonder ijsdek}}}{C_{\text{met ijsdek}}}\right)^{\frac{2}{3}}\approx \left(\frac{1}{\sqrt{2}}\right)^{\frac{2}{3}},
\end{equation}
waarbij $C$ de Chezy ruwheidscoëfficient is [$m^{\frac{1}{2}}/s$].

Het Froude getal $Fr$ is gedefinieerd als
\begin{equation}
	Fr = {\frac{u}{\sqrt{gh}}} = {\frac{q}{h\sqrt{gh}}},
\end{equation}
waarbij $u$ de stroomsnelheid is [$m/s$], $q$ het specifieke debiet [$m^2/s$] en $g$ de gravitatie-constante [$m/s^2$]. Het effect van de opstuwing op $Fr$ wordt dan geschat met
\begin{equation}
	Fr_{\text{met ijsdek}} \approx \left({\frac{h_{\text{zonder ijsdek}}}{h_{\text{met ijsdek}}}}\right)^{\frac{3}{2}}
\end{equation}
\begin{equation}
	Fr_{\text{zonder ijsdek}} = {\frac{C_{\text{met ijsdek}}}{C_{\text{zonder ijsdek}}}}
\end{equation}
\begin{equation} \label{eq:5}
	Fr_{\text{met ijsdek}} = {\frac{Fr_{\text{zonder ijsdek}}}{\sqrt{2}}}
\end{equation}
De correctie voor opstuwing door een benedenstrooms ijsdek volgt dus eenvoudig uit $0.71Fr$, berekend in D-HYDRO. Hier is verondersteld dat stroombanen in de hoofdgeul onder invloed van het benedenstrooms ijsdek niet verschuiven. Door deze vereenvoudiging kan $Fr$ wat worden overschat omdat bij opstuwing in de hoofgeul de uiterwaarden wat gemakkelijker meestromen. 

\textbf{Correctie 2}\\
De invloed van lokale bodemveranderingen op het Froude getal wordt vergelijkbaar geschat met
\begin{equation} \label{eq:6}
	\frac{Fr_{z1}}{Fr_{z0}} \approx \left({\frac{h_{z0}}{h_{z1}}}\right)^{\frac{3}{2}} = \left(1-{\frac{\Delta{z}}{h_{z0}}}\right)^{-\frac{3}{2}},
\end{equation}
waarbij $z$ de bodemligging is. $z1$ geeft bodemverandering aan en $z0$ géén bodemverandering.

\textbf{Gecombineerde correctie}\\
De uiteindelijke, gecombineerde correctie van het met D-HYDRO berekende Froude getal is dan (\ref{eq:5}, \ref{eq:6})  
\begin{equation}
	\frac{Fr_{gecorrigeerd}}{Fr_{berekend}} = \frac{1}{\sqrt{2}}\left(1-{\frac{\Delta{z}}{h_{z0}}}\right)^{-\frac{3}{2}}.
\end{equation}

\fi

\LastPage
\end{document}
